\subsection{Source and Receiver Node Architecture}

For repeatability and control of the experiment setup, the system employs two separate embedded nodes: a fixed transmitter node and an autonomous receiver node with a mechanically steerable directional antenna. Functional roles of the two nodes are carefully segregated to decouple antenna orientation from the received signal strength.

\subsubsection{Transmitter Node}

A transmitter node consists of an ESP32-S3 development board with a directional antenna tuned to the 2.4~GHz ISM band. Antenna orientation of the transmitter node is fixed throughout all experiments. As a result, a static but non-uniform radiation pattern in space is created, which leads to an RSSI distribution that depends on both the radiation pattern of the transmitter antenna and the orientation of the receiver antenna. Such a setup reflects more realistic point-to-point wireless communication systems when both sides have directional antennas to optimize link budget and minimize interference.

The transmitter node emits wireless packets at a steady rate and with a constant power level. No adaptive control of the antenna is used in this experiment scenario, which means any variation of RSSI measurements observed at the receiver side is purely due to antenna misalignment, channel properties, and receiver control activity.

\subsubsection{Receiver Node}

The receiver node consists of an ESP32-S3 microcontroller combined with a high-gain directional antenna capable of mechanical steering along two axes: azimuth and elevation. Both of these directions allow fine-tuning antenna orientation to maximize RSSI reception towards the transmitter node. On receiving a wireless packet, the ESP32-S3 obtains RSSI information and uses it as a feedback for further antenna adjustment. To attenuate fading and noise, RSSI values can be smoothed using simple averaging techniques before any decision.

All alignment actions are calculated in the receiver node. Depending on the selected control policy, the ESP32-S3 controller computes movement commands for pan/tilt actuator drivers that control the antenna position via one-way communication.